\section{Additional experiments and analysis} \label{app:more-experiments}


\paragraph{Model variance}
The main experiments in Section~\ref{sec:experiments} are based on a single run.
We ran automated evaluation on 3 random runs of \approach{}, using PaLM outputs on NQ as input passages.
The standard deviations of $\attrauto{}$, $\preslev{}$, and $\HM{}$ are 1.2, 0.5, and 1.0 respectively.

\paragraph{Impact of the retriever choice}
We tried using Microsoft Bing in place of Google Search, with near identical results (< 1\% difference).

\begin{table*}[!t]
\centering\small\setlength{\tabcolsep}{5pt}
\adjustbox{max width=\textwidth}{%
\begin{tabular}{
@{}l
r@{ $\to$ }rcc
r@{ $\to$ }rcc
r@{ $\to$ }rcc
@{}}
\toprule
 & \multicolumn{4}{c}{PaLM outputs on NQ}
 & \multicolumn{4}{c}{PaLM outputs on SQA}
 & \multicolumn{4}{c}{LaMDA outputs on QReCC}
 \\
 \cmidrule(lr){2-5} \cmidrule(lr){6-9} \cmidrule(lr){10-13}
 Model
 & \multicolumn{2}{c}{$\attrauto{}$} & $\preslev{}$ & \multicolumn{1}{c}{$\HM{}$}
 & \multicolumn{2}{c}{$\attrauto{}$} & $\preslev{}$ & \multicolumn{1}{c}{$\HM{}$}
 & \multicolumn{2}{c}{$\attrauto{}$} & $\preslev{}$ & \multicolumn{1}{c@{}}{$\HM{}$}
 \\
\midrule
Full \approach{}
& 45.6 & 54.9 & 89.6 & 68.1
& 37.6 & 45.1 & 89.9 & 60.0
& 18.8 & 29.4 & 80.2 & 43.1
\\ %
qgen 62B, editor 540B
& 45.9 & 54.6 & 87.8 & 67.4
& 37.0 & 40.5 & 90.0 & 55.9
& 15.8 & 28.4 & 76.1 & 41.4
\\
qgen 62B, editor 62B
& 45.9 & 49.9 & 91.0 & 64.4
& 37.0 & 38.3 & 93.0 & 54.2
& 15.8 & 21.9 & 71.6 & 33.5
\\ %
GPT-3
& 44.3 & 55.0 & 90.6 & \textbf{68.5}
& 38.6 & 46.6 & 89.3 & \textbf{61.2}
& 18.3 & 28.6 & 89.8 & \textbf{43.4}
\\
\bottomrule
\end{tabular}
}
\caption{\textbf{Additional ablation results.} We report the automatic metrics: $\attrauto$, $\preslev$, and harmonic mean between the two ($\HM{}$). We show auto-AIS scores both before and after editing (before $\rightarrow$ edit), with respect to the attribution report $A$ produced by the model.
}
\label{tab:ablation_results_upper}
\end{table*}

\paragraph{Impact of model scale}
Many components in \approach{} work by few-shot prompting PaLM, a large 540B  parameter LM. To assess the benefit of LM scaling, we replaced PaLM 540B with a smaller 62B parameter PaLM. As shown in Table~\ref{tab:ablation_results_upper}, we found that 540B outperforms 62B by a large margin, suggesting that \approach{} could potentially further improve with even more scaling. We also experimented with keeping the editor stage at 540B while shrinking the query generation stage to 64B --- this yielded a relatively small performance drop, suggesting that model scaling is more important for the editor.

\paragraph{Impact of model type}
Few-shot prompting has proven to be effective for many recent large language models. We try replacing the query generation model, agreement model, and edit model with GPT-3 text-davinci-003. The few-shot prompts were slightly tuned to fit the GPT-3 model.
Table~\ref{tab:ablation_results_upper} shows the results, which are slightly better than \approach{} implemented with PaLM 540B on all three datasets.
We will release this open-source version of RARR that uses GPT-3 as the backbone.








\begin{table*}[!t]
\centering\small
\adjustbox{max width=\textwidth}{%
\begin{tabular}{
@{}l
r@{ $\to$ }rcc
r@{ $\to$ }rcc
r@{ $\to$ }rcc
@{}}
\toprule
 & \multicolumn{4}{c}{GPT-3 outputs on NQ}
 & \multicolumn{4}{c}{GPT-3 outputs on SQA}
 & \multicolumn{4}{c}{GPT-3 outputs on QReCC}
 \\
 \cmidrule(lr){2-5} \cmidrule(lr){6-9} \cmidrule(lr){10-13}
 Model
 & \multicolumn{2}{c}{$\attrauto{}$} & $\preslev{}$ & \multicolumn{1}{c}{$\HM{}$}
 & \multicolumn{2}{c}{$\attrauto{}$} & $\preslev{}$ & \multicolumn{1}{c}{$\HM{}$}
 & \multicolumn{2}{c}{$\attrauto{}$} & $\preslev{}$ & \multicolumn{1}{c@{}}{$\HM{}$}
 \\
\midrule
EFEC
& 48.3 & 66.8 & 41.5 & 51.2
& 32.6 & 50.6 & 29.4 & 37.2
& 26.4 & 53.1 & 39.0 & 44.9
\\
LaMDA
& 36.2 & 61.1 & 45.9 & 52.4
& 22.3 & 27.3 & 43.3 & 33.5
& 19.0 & 33.9 & 28.3 & 30.8
\\
PaLM \approach{}
& 48.3 & 57.2 & 89.6 & 69.8
& 32.6 & 36.3 & 91.6 & 52.0
& 26.4 & 31.1 & 87.7 & \textbf{45.9}
\\
GPT-3 \approach{}
& 48.0 & 59.3 & 91.8 & \textbf{72.0}
& 34.7 & 37.0 & 91.8 & \textbf{52.8}
& 23.2 & 25.3 & 89.7 & 39.5
\\
\bottomrule
\end{tabular}
}
\caption{\textbf{Results on passages from GPT-3.} We report the automatic metrics: $\attrauto$, $\preslev$, and harmonic mean between the two ($\HM{}$). We show auto-AIS scores both before and after editing (before $\rightarrow$ edit), with respect to the attribution report $A$ produced by the model.
The results show similar trends as the results on passages from PaLM and LaMDA in Table~\ref{tab:main_results}.
}
\label{tab:gpt3_passages_results}
\end{table*}

\paragraph{Results on GPT-3 passages}
Table~\ref{tab:gpt3_passages_results} shows automated evaluation results on passages generated by GPT-3.
The results follow the same trend as the results on PaLM and LaMDA passages.

\paragraph{Challenging domains}
We report results on tasks where attribution was particularly hard, and significant future work is needed.

\begin{table}[!t]
\centering\small
\adjustbox{max width=\columnwidth}{%
\begin{tabular}{
@{}l
r@{ $\to$ }rcc
@{}}
\toprule
 Model
 & \multicolumn{2}{c}{$\attrauto{}$} & $\preslev{}$ & $\HM{}$
 \\ \midrule
\multicolumn{5}{c}{\textbf{SummEval}}
\\
EFEC 
& 17.9 & 34.6 & 20.9 & 26.0 
\\
LaMDA
& 10.3 & 28.8 & 28.1 & 28.4 
\\
\approach{}
& 18.3 & 16.9 & 92.9 & 28.6 
\\
\midrule
\multicolumn{5}{c}{\textbf{ELI5}}
\\
EFEC 
& 18.2 & 41.2 & 17.2 & 24.2 
\\
LaMDA
& 19.9 & 40.1 & 31.2 & 35.1 
\\
\approach{}
& 18.5 & 18.9 & 97.2 & 31.7
\\

\bottomrule
\end{tabular}
}
\caption{\textbf{Results on ELI5 and SummEval.}}
\label{tab:eli5_and_summeval}
\end{table}



We considered news article summaries produced by summarization models from SummEval \citep{fabbri21summeval} (e.g., \textit{``John Doe was left homeless when the storms hit Staten Island, New York \dots''}). Results are shown in Table~\ref{tab:eli5_and_summeval}. First, we note that the before-edit auto-AIS scores for all models are low. These news article summaries are often about less widely known people and events, which is challenging for retrievers, leading to low attribution. For example, our query generator may ask \emph{``where does John Doe live''} but get results for a different John Doe. EFEC and LaMDA also face this issue, but instead trade preservation for attribution and rewrite the text to a different topic. 
This result suggests that using web search with standard question generation methods may fail to capture important context from the input, and is not sufficient for the attribution task.

We also considered long-form explanations generated by PaLM for the ELI5 dataset~\citep{fan19eli5} (Table~\ref{tab:eli5_and_summeval}). ELI5 was collected from online forums, so many answers tend to have subjective opinions instead of specific entities and facts (e.g., \textit{``How do our brains interpret scary music? To me, scary music often sounds a little bit like a person \dots''}), and are thus difficult to attribute. Sometimes the whole output is based on a false premise and needs to be completely rewritten, in which case \approach{} cannot satisfactorily edit due to our revision threshold (Section~\ref{sec:approach-editor}). %

\begin{table}[!t]
\centering\small
\adjustbox{max width=\columnwidth}{%
\begin{tabular}{
@{}l
r@{ $\to$ }rcc
@{}}
\toprule
 & \multicolumn{4}{c}{RARR}
 \\
MMLU Category
 & \multicolumn{2}{c}{$\attrauto{}$} & $\preslev{}$ & $\HM{}$
 \\ \midrule
\text{Humanities} 
& 26.6 & 29.6 & 6.6 & 45.0 
\\
\text{Social Sciences}
& 35.5 & 40.7 & 7.6 & 56.5
\\
\text{STEM}
& 37.8 & 41.5 & 7.2 & 57.4 
\\
\text{Other}
& 36.9 & 41.7 & 7.1 & 57.6 
\\ 
\bottomrule
\end{tabular}
}
\caption{\textbf{\approach{} results on MMLU.}}
\label{tab:mmlu}
\end{table}



Finally, we considered technical explanations to questions from the MMLU dataset \cite{hendryckstest2021} which covers diverse subjects from social science, humanities, STEM, and others.\footnote{MMLU has questions from 57 subjects; we took 10 random question from each topic and generated answer explanations by prompting PALM 540B.} An example input looks like \textit{ ``Every time you remove an edge from a complete graph, you divide it into two connected components. So, a complete graph with 13 vertices must have 12 connected components.''} Results are shown in Table~\ref{tab:mmlu}. \approach{} improves attribution of the explanations on all four categories of MMLU, although the increases are relatively small. We also found that \approach{}'s performance is low on examples with mathematical reasoning, as these are beyond the capability of the edit model with our current prompt.

\section{Details on automated evaluation}\label{sec:app-auto-eval}
\paragraph{Sentence splitting} When computing the attribution score, we use spaCy \verb|en_core_web_sm| v3.0.0a1 to segment the text passage into sentences. (More recent models gave similar results.) While each sentence may contain multiple claims that could be attributed independently, there is currently no linguistic consensus on what constitutes a claim. Instead of depending on a particular definition of claims, we use sentences as claims for simplicity and reproducibility. The same segmentation is also used for human evaluation.

\paragraph{Decontextualization} We decontextualize each sentence in the text passage before computing the attribution score. We use the model from \citet{choi2021making}, which is a T5 model fine-tuned to map the input ``\texttt{[HEAD] [SEP] }\textit{context and passage}\texttt{ [start] }\textit{sentence}\texttt{ [end]}'' to the output ``\texttt{[OPCODE] }\textit{decontextualized sentence}'', where the \texttt{OPCODE} can be ``\texttt{done}'' (success), ``\texttt{un}'' (unnecessary), or ``\texttt{imp}'' (impossible). We feed the passage's context (questions for NQ and SQA; dialog context for QRECC) along with the passage itself to the input. We use beam search with beam size 8 and discard any result whose number of tokens differ by more than 4.

\paragraph{NLI model} We obtained a newer version of the end-to-end NLI model from the authors of \citet{Honovich2022TRUERF}, which was trained on MNLI, SNLI, FEVER, PAWS, SciTail and VitaminC \cite{williams2017broad,bowman2015large,thorne2018fever,zhang2019paws,khot2018scitail,schuster2021get}.
The model is a T5 model fine-tuned to map the input ``\texttt{premise: }\textit{evidence}\texttt{ hypothesis:   }\textit{claim sentence}'' to either ``\texttt{1}'' (entailed) or ``\texttt{0}'' (not entailed). As suggested by the authors, we use the probability of producing ``\texttt{1}'' as the entailment score.

\begin{figure}[!t]
\centering
\includegraphics[width=\columnwidth]{figures/human-vs-auto-eval2.png}
\caption{
Violin plot illustrating the strong correlation between human AIS and auto-AIS labels on our NQ benchmark. Pearson correlation is 0.74 (N=450). $y$-axis is auto-AIS score, the two violins correspond to a human label of 0 or 1.
}\label{fig:human-vs-auto-eval}
\end{figure}

\paragraph{Comparing human and automated evaluation}
We conducted correlation studies between human and automatic metrics and found strong Pearson correlation (attribution = 0.74; preservation = 0.62).
We visualize the correlation between human and automated attribution scores on NQ and SQA in Figure~\ref{fig:human-vs-auto-eval}. We found that the AIS scores from human correlate well with auto-AIS scores, with some bias for non-attributed sentences to be judged as attributed by auto-AIS.

\section{Details on human evaluation}
\label{sec:app-human-eval}

To end-goal of \approach{} is to improve the \textit{attribution} of generation models through post-editing while \textit{preserving} the original intent.
Attribution and preservation are both subjective properties that may change with even small edits.
In the main paper, we present two automatic metrics to conveniently gauge these properties, but rely on a human evaluation as the gold standard.
In this section, we describe how we conducted the human evaluation and what instructions and examples annotators were provided.

\paragraph{Rater recruitment and training}
We engaged with a vendor supplier of full-time crowd workers to recruit human annotators for our task.
Annotators were asked to review the instructions below and were provided direct feedback on their responses during the pilot annotation runs.
We had 3 annotators rate each example in the pilot phase to measure inter-annotator agreement, and had a single rater annotate each example afterwards.

\subsection{Instructions: Overview}

In this task you will evaluate the quality of text generated by a system (\textbf{the ``passage''}) based on how well it represents information from multiple pieces of \textbf{``evidence''}.

We will be using two categories to evaluate the quality of the passage: \textbf{Attribution} and \textbf{Intent Similarity}. You will evaluate these categories in succession. In some tasks, you will only evaluate Attribution. The task interface will guide you through the flow; you can also see the overall task flow in the diagram below. 

\textbf{Note:} The passage may appear very fluent and well-formed, but still contain slight inaccuracies that are not easy to discern at first glance. Pay close attention to the text. Read it carefully as you would when proofreading.

\subsection{Instructions: Attribution}

\newcommand{\attrspan}[1]{\sethlcolor{lime}\hl{#1}}
\newcommand{\nattrspan}[1]{\sethlcolor{pink}\hl{#1}}

\begin{figure*}[!thp]
    \centering
    \includegraphics[width=\textwidth]{figures/human-eval-attribution}
    \caption{
    Screenshot of interface to annotate attribution at the sentence level.
    annotators were asked to mark sentences as being \attrspan{fully attributable} or \nattrspan{not fully attributable} by clicking each sentence, and rating each piece of evidence as being useful or not in helping determine attribution of the passage. Annotators were also presented with the context of the generation.
    }
    \label{fig:human-eval-attribution}
\end{figure*}
\begin{figure*}[!thp]
    \centering
    \includegraphics[width=\textwidth]{figures/human-eval-preservation.png}
    \caption{Screenshot of the preservation interface. Annotators are asked to read compare two passages and rate how similar the intent conveyed by the two passages is.}
    \label{fig:human-eval-preservation}
\end{figure*}

In this step, you will evaluate how much of the passage is attributable to one or more pieces of evidence (Figure \ref{fig:human-eval-attribution}).

In the interface, the passage of text and the context in which it was generated is shown on the left, and each piece of evidence is shown on the right. You will use all three (context, passage, evidence) to answer the following question for each sentence in the passage:
\textit{Is all of the information provided by this sentence fully supported by at least one piece of evidence?}

\paragraph{Determining the information provided by the sentence.}
Three points are key when determining information provided by the sentence:
\begin{enumerate}
\item The context and the other sentences of the passage are often critical in understanding the information provided by the sentence.
\item The context should only be used to understand the information provided by the sentence.
\item The evidence should be completely ignored for this step.
\end{enumerate}

Consider the following example:

\begin{quote}
\textit{Context:} who plays doug williams in days of our lives \\
\textit{Passage:} In the American daytime drama Days of Our Lives, Doug Williams and Julie Williams are portrayed by Bill Hayes and Susan Seaforth Hayes. \\
\end{quote}

In the above example, the meaning of the passage is clear even without seeing the query. But consider another example:

\begin{quote}
\textit{Context:} who plays doug williams in days of our lives \\
\textit{Passage:} he is played by Bill Hayes \\
\textit{Passage (interpreted):} \ul{Doug Williams} is played by Bill Hayes \ul{in days of our lives}
\end{quote}

In this case the pronoun ``he'' depends on the context, but it is clear that the  intended meaning of the passage can be reasonably interpreted as ``Doug Williams is played by Bill Hayes in days of our lives''. This interpretation is the ``information provided by the passage''. 

Pronouns such as he/she/it/they etc. are one case where context is needed to figure out the intended meaning of the system response. Here's another example (given with paraphrases of the information highlighted below):

\begin{quote}
\textit{Context:} when is the last time the us lost basketball at the olympics \\
\textit{Passage:} \ul{The last time they lost was in 2004}, when Argentina defeated the US 89–79. Most recently, they won gold in 2016. \\
\textit{Passage (interpreted):} The last time \ul{the United States lost basketball} at the Olympics was in 2004.
\end{quote}

The context should only be used to determine the information provided by the passage; at times, the passage may be about a slightly different topic than the context, for example:

\begin{quote}
\textit{Context:} the south west wind blows across nigeria between \\
\textit{Passage:} \ul{The Harmattan is a dry and dusty northeasterly trade wind that blows across West Africa} from December to March. It is very dusty because it blows across the Sahara.
\end{quote}

Here, the passage talks about a \textit{northeasterly wind}, while the context asks about a \textit{south-west wind}, but the passage can be fully understood.

In general, use your best judgment to determine the information provided by the passage. If the passage is hard to understand and you are unsure what the intended meaning of the passage is, \textit{mark the sentences as not attributed} and enter a comment with an explanation. As one example, take the following:

\begin{quote}
\textit{Context:} how many NBA championships did Michael Jordan win? \\
\textit{Passage:} it is the best team in the NBA
\end{quote}

\paragraph{Determining if the information accurately represents the evidence.}
Two points are key when determining whether the information accurately represents the evidence:
When interpreting a piece of evidence, use only the title and text of that specific evidence. Completely ignore the context, passage and all other evidence.
Check all the information in a sentence. If only some information is supported by the evidence, mark the sentence as not fully attributable.

Consider the following example:
\begin{quote}
\textit{Context:}
when did reba mcentire record back to god \\
\textit{Passage:}
\attrspan{Back to God was released by McEntire in 2017.} \\
\textit{Evidence:}
\ul{``Back to God'' is a song performed by American singer, Reba McEntire.} \ul{It was released} as the second single from her 2017 album, Sing it Now: Songs of Faith \& Hope, \ul{on January 20, 2017}.
\end{quote}
In the above example, it is reasonable to conclude that the evidence supports all the information in the passage, and we can mark the passage as being fully attributable. But consider another example:

\begin{quote}
\textit{Context:}
who won the womens 2017 ncaa basketball tournament \\
\textit{Passage:}
\nattrspan{South Carolina Gamecocks won the 2017 NCAA Women's Division I Basketball Tournament.} \\
\textit{Evidence:}
The \ul{South Carolina Gamecocks} defeated the Mississippi State Bulldogs, 67–55, to claim their first-ever \ul{national championship}.
\end{quote}

In this case, while the evidence also mentions the ``South Carolina Gamecocks'', it isn't clear that the national championship being mentioned is indeed the 2017 NCAA Women's Division I Basketball Tournament. The passage should be marked as not attributable.

Finally, when the passage contains multiple sentences, evaluate whether each sentence can be fully attributed to one or more pieces of evidence—it is possible for one sentence to be attributed while another is not. For example:

\begin{quote}
\textit{Context:}
who won the womens 2017 ncaa basketball tournament \\
\textit{Passage:}
\nattrspan{South Carolina Gamecocks won the 2017 NCAA Women's Division I Basketball Tournament.} \nattrspan{The final score is 67-55.} \attrspan{The championship game was held in Dallas, Texas.} \\
\textit{Evidence 1:}
The South Carolina Gamecocks defeated the Mississippi State Bulldogs, 67–55, to claim their first-ever national championship. \\
\textit{Evidence 2:}
The \ul{2017 NCAA Women's Division I Basketball Tournament} was played from Friday, March 17 to Sunday, April 2, 2017, with the Final Four played at the American Airlines Center in \ul{Dallas, Texas} on March 31 and April 2.
\end{quote}

The first two sentences cannot be attributed to either evidence for the same reason as the previous example, but the last sentence is fully supported by Evidence 2 and should be marked as attributed.

In general, you should use your best judgment in determining whether all of the information provided by the passage is ``an accurate representation of information in at least one evidence''. See Table \ref{tab:human-eval-attribution-examples} for additional examples.

We give the following final notes of guidance:
\begin{itemize}
\item  \textbf{Marking evidence as useful.} When reviewing each piece of evidence, mark it as useful if it helps you judge the attributability of any sentence, and mark it not useful if not. In the above example Evidence 1 is not useful because it didn't contain enough context to actually help you assess if the passage was attributable, but Evidence 2 was useful.
\item
\textbf{Contradicting evidence.} Mark a sentence as being attributed if any piece of evidence supports it: if two pieces of evidence contradict each other, but one of them supports the passage, mark the sentence as fully attributable.
\item
\textbf{More on the concept of ``accurate representation''.} We take as inspiration the journalist's conception of ``accurate representation''. For example, take \href{https://www.npr.org/about-npr/688139552/accuracy}{this excerpt on Accuracy in the NPR Ethics Handbook}: ``When quoting or paraphrasing anyone … consider whether the source would agree with the interpretation...'' In other words, if you had written the source document, consider whether you would view the system response as an accurate representation of information in that source document.
\end{itemize}


\begin{table*}
    \centering \small
    \begin{tabular}{@{}p{.25\textwidth} p{.45\textwidth} p{.25\textwidth}@{}}
    \toprule
        Context + Passage &
        Evidences &
        Notes \\ \midrule
        
\textit{Context: who played morticia in the addams family tv show}

\attrspan{The Addams Family is an American animated sitcom TV series.}
\nattrspan{It was first aired on NBC on September 24, 1973. Carolyn Jones played the role of Morticia.}
&
1/ The Addams Family (1973 TV series): The Addams Family is an American animated sitcom adaptation of the Charles Addams comic. The series was produced in 1973 and was rebroadcast the following season.

2/ The Addams Family (TV Series 1964–1966): When The Addams Family went off the air in 1966, network executives in charge of children's programming for NBC brought them back in 1973 for their own Saturday Morning cartoon show featuring the voices of Carolyn Jones from the original series.
&
While the evidence supports the show being aired in 1973, it doesn't specify the exact date (September 24, 1973).

Similarly, while the evidence mentions Carolyn Jones as being a voice actor, it doesn't say she played the role of Mortica.
\\      \midrule
\textit{
Context: when will the la sagrada familia be finished
}

\attrspan{The La Sagrada Familia is a large Roman Catholic church in Barcelona. It is designed by Antoni Gaudi.}
\nattrspan{It started construction in 1882, and the construction is still going on.}
\nattrspan{The estimated date to finish is 2026.}
&

1/ Sagrada Família - Wikipedia: The Basílica i Temple Expiatori de la Sagrada Família is a church in the Eixample district of Barcelona, Catalonia, Spain, and is currently the
largest unfinished Roman Catholic church.

2/ Find Out Sagrada Familia's Expected Finish Date: Visiting the breathtaking Sagrada Familia today also means witnessing the slow progress towards the completion of the project. Sagrada Familia is now  expected to be completed in 2026, the centenary of Gaudi's death.
It's a reasonable inference that La Sagrada Familia is the same as Sagrada Familia, even though the names differ slightly.
&
While Evidence 2 mentions Gaudi, it isn't clear this is a reference to Antoni Gaudi and further doesn't say that he designed the church.
\\
    \bottomrule
    \end{tabular}
    \caption{Additional examples for annotating attribution.}
    \label{tab:human-eval-attribution-examples}
\end{table*}

\subsection{Instructions: Intent Similarity}
In this step, you will evaluate how much similar the passage is to another passage (Figure \ref{fig:human-eval-preservation}).

In the interface, the passage A and passage B are both text generated by a system—given the same context in which it was generated. You will use all three (context, passage A, passage B) to answer the following question:
\textit{
How similar is the intent expressed by Passage A and Passage B? Please ignore any differences in details.
}

Two points are key when determining whether the two passages convey the same intent:
\begin{enumerate}
    \item  Judge the similarity solely based on the similarity in the type and quantity of information provided by each passage.
    \item Ignore any differences in factual details between the two passages.
\end{enumerate}

Consider the following examples:

\begin{quote}
\textit{Context:}
who pays medical bills in great britain where does the money come from to pay these bills \\
\textit{Passage A:}
Britain's National Health Service (NHS) is paid for through general taxation \ul{and national insurance}. In 2017/18, the NHS budget was £\ul{176.5bn}. \\
\textit{Passage B:}
Britain's National Health Service (NHS) is paid for through general taxation. In 2017/18, the NHS budget was £\ul{118bn}. \\
\textit{Rating:} \ul{Very similar.} Passage A is about the same topic as Passage B, with a similar level of detail and style of presentation. They may differ in factual details.
\end{quote}

The above example should be rated ``very similar'' because both passages include information about (1) how the NHS is paid for, and (2) what its budget in 2017/18 was, though they differ in their actual answers to these questions.

\begin{quote}
\textit{Context:}
who is the owner of reading football club \\
\textit{Passage A:}
Reading's owner is \ul{Yongge Dai}. Yongge \ul{Dai is also the president of Chinese company Dai Yongge Real Estate. Yongge's son, Dai Xiu Li, is Reading's vice-president.} \\
\textit{Passage B:}
Reading's owner is \ul{Dai Yongge}. Yongge's \ul{brother and sister pair behind the Reading FC takeover---Dai Yongge and Dai Xiu Li---has made their fortune through a massive property empire. Mr Dai, has been the chairman of Renhe Commercial since 1999, which is an organisation owned by his sister behind a vast network of underground shopping centres in China}. \\
\textit{Rating:}
\ul{Somewhat similar.} Passage A is about the same topic as Passage B, but differs substantially in level of detail or style of presentation. They may differ in factual details.
\end{quote}

The above example should be rated ``somewhat similar'' because both passages are still about the same topic—Reading's owner— but differ substantially in the information they discuss: Passage A includes information about (1a) who Reading's owner is, (2a) which company they are the president of and (3a) who their vice-president is. In contrast, while Passage B shares information about (1a), it also includes information about (2b) how the Reading owner made their fortune, (3b) their company position and how long they held it for and (4b) what the company also owns.


\begin{quote}
\textit{Context:}
what is the numbers of total elected member of indian parliment in present time \\
\textit{Passage A:}
The total number of elected members of the \ul{Lok} Sabha is \ul{543}. \\
\textit{Passage B:}
The total number of elected members of the \ul{Rajya} Sabha is \ul{238}. \\
\textit{Rating:}
\ul{Not at all similar.} Passage A is about a significantly different topic than Passage B.
\end{quote}

Even though the passages look very similar, the above example should be rated ``not at all similar'' because the two passages are about significantly different topics: ``the Lok Sabha'' vs ``the Rajya Sabha''.


\section{Details on the model}
\label{app:model-details}

\paragraph{Few-shot prompting with LLMs} \label{section:prompting}

We implement many sub-tasks within \approach{} using \emph{few-shot prompting} of LLMs (also known as \emph{in-context learning}~\cite{Brown2020LanguageMA}) as follows:
\begin{enumerate} \setcounter{enumi}{0}
\item
For each sub-task, we manually author a small number of training examples: \texttt{$(\text{input}_j, \text{output}_j)$} for $j=1,\ldots,J$, where $J$ ranges between 5 and 10 and where both the input and output are strings.
\item
We form the following prompt: \texttt{$\text{input}_1 \diamond \text{output}_1 \oplus \text{input}_2 \diamond \text{output}_2 \oplus \ldots \oplus \text{input}_J \diamond \text{output}_J \oplus \text{new\_input}$}, where $\diamond$ denotes a newline character and $\oplus$ denotes a double newline character.
\item
To perform inference on a new input, we condition the LLM on the prompt and sample continuations of the prompt up until the next double newline character.
\end{enumerate}
All of our prompts are included in Figures~\ref{fig:prompt_cqgen}, \ref{fig:prompt_agreement}, and~\ref{fig:prompt_revision}. The contextual version used for QReCC are in Figures~\ref{fig:prompt_cqgen_contextual}, \ref{fig:prompt_agreement_contextual}, and~\ref{fig:prompt_revision_contextual}.


\paragraph{Model statistics}


We implemented most parts of RARR with the PALM model which has 540B parameters. We prompted PALM without any training or finetuning. We used a TPU  v2-128 to run inference with PALM.

We manually wrote our prompts by eye-balling quality on a dozen of examples from a separate validation set.
We tune our hyperparameters on the validation set as well. We used sampling temperature 0.7 for all generation tasks. For each input text, we sample 3 question generations, and for each question we retrieve 5 results. For agreement gate and editing, we only sample 1 generation. We reject an editing if the edit distance is more than 50 characters or more than half of the original text length.

\section{Details on the dataset}

As explained in Section~\ref{sec:eval-setup}, we generated 150 development and 150 test passages for each of the 6 combinations of dataset and model: (NQ, PaLM), (SQA, PaLM), (QReCC, LaMDA), (NQ, GPT-3), (SQA, GPT-3), (QReCC, GPT-3).
Figures~\ref{fig:prompt_passage_palm}, \ref{fig:prompt_passage_gpt3}, \ref{fig:prompt_passage_lamda},  and~\ref{fig:prompt_passage_gpt3_conv} are the few-shot prompts used to generate the passages.

Following the corresponding datasets, all generated passages are in English. The authors have manually looked through most of the data and found no personal identifiers.




{ %

\newcommand{\pWidth}{5.8in}
\newcounter{pCount}
\newcommand{\pStart}{\setcounter{pCount}{1}}
\newcommand{\pRow}[1]{\texttt{\tiny\color{gray}\arabic{pCount}}&\texttt{#1}\stepcounter{pCount}\\}
\newcommand{\pVar}[1]{\textbf{\{#1\}}}
\newcommand{\pBlank}{\_\_\_\_\_}

\begin{figure*}[p]
\begin{center}\scriptsize
\begin{tabular}{|rp{\pWidth}|}\hline\pStart&\\
\pRow{[web] I will check things you said and ask questions.}
\pRow{}
\pRow{(1) You said: Your nose switches back and forth between nostrils. When you sleep, you switch about every 45 minutes.  This is to prevent a buildup of mucus. It's called the nasal cycle.}
\pRow{To verify it,}
\pRow{a) I googled: Does your nose switch between nostrils?}
\pRow{b) I googled: How often does your nostrils switch?}
\pRow{c) I googled: Why does your nostril switch?}
\pRow{d) I googled: What is nasal cycle?}
\pRow{}
\pRow{(2) You said: The Stanford Prison Experiment was conducted in the basement of Encina Hall, Stanford's psychology building.}
\pRow{To verify it,}
\pRow{a) I googled: Where was Stanford Prison Experiment was conducted?}
\pRow{}
\pRow{(3) You said: The Havel-Hakimi algorithm is an algorithm for converting the adjacency matrix of a graph into its adjacency list. It is named after Vaclav Havel and Samih Hakimi.}
\pRow{To verify it,}
\pRow{a) I googled: What does Havel-Hakimi algorithm do?}
\pRow{b) I googled: Who are Havel-Hakimi algorithm named after?}
\pRow{}
\pRow{(4) You said: "Time of My Life" is a song by American singer-songwriter Bill Medley from the soundtrack of the 1987 film Dirty Dancing. The song was produced by Michael Lloyd.}
\pRow{To verify it,}
\pRow{a) I googled: Who sings "Time of My Life"?}
\pRow{b) I googled: Which film is "Time of My Life" from?}
\pRow{c) I googled: Who produced the song "Time of My Life"?}
\pRow{}
\pRow{(5) You said: Kelvin Hopins was suspended from the Labor Party due to his membership in the Conservative Party.}
\pRow{To verify it,}
\pRow{a) I googled: Why was Kelvin Hopins suspended from Labor Party?}
\pRow{}
\pRow{(6) You said: Social work is a profession that is based in the philosophical tradition of humanism. It is an intellectual discipline that has its roots in the 1800s.}
\pRow{To verify it,}
\pRow{a) I googled: What philosophical tradition is social work based on?}
\pRow{b) I googled: What year does social work has its root in?}
\pRow{}
\pRow{(7) You said: \pVar{text}}
\pRow{To verify it,}
\pRow{\pBlank}
&\\\hline\end{tabular}
\caption{Few-shot prompt for query generation. To increase diversity and coverage, we sample the model three times and combine the resulting lists of queries.}
\label{fig:prompt_cqgen}
\end{center}
\end{figure*}

\begin{figure*}[p]
\begin{center}\scriptsize
\begin{tabular}{|rp{\pWidth}|}\hline\pStart&\\
\pRow{[web] I will check some things you said.}
\pRow{}
\pRow{(1) You said: Your nose switches back and forth between nostrils. When you sleep, you switch about every 45 minutes.  This is to prevent a buildup of mucus. It's called the nasal cycle.}
\pRow{I checked: How often do your nostrils switch?}
\pRow{I found this article: Although we don't usually notice it, during the nasal cycle one nostril becomes congested and thus contributes less to airflow, while the other becomes decongested. On average, the congestion pattern switches about every 2 hours, according to a small 2016 study published in the journal PLOS One.}
\pRow{Your nose's switching time is about every 2 hours, not 45 minutes.}
\pRow{This disagrees with what you said. }
\pRow{}
\pRow{(2) You said: The Little House books were written by Laura Ingalls Wilder. The books were published by HarperCollins.}
\pRow{I checked: Who published the Little House books?}
\pRow{I found this article: These are the books that started it all -- the stories that captured the hearts and imaginations of children and young adults worldwide. Written by Laura Ingalls Wilder and published by HarperCollins, these beloved books remain a favorite to this day.}
\pRow{The Little House books were published by HarperCollins.}
\pRow{This agrees with what you said.}
\pRow{}
\pRow{(3) You said: The Stanford Prison Experiment was conducted in the basement of Jordan Hall, Stanford's psychology building.}
\pRow{I checked: Where was Stanford Prison Experiment conducted?}
\pRow{I found this article: Carried out August 15-21, 1971 in the basement of Jordan Hall, the Stanford Prison Experiment set out to examine the psychological effects of authority and powerlessness in a prison environment.}
\pRow{The Stanford Prison Experiment was conducted in Jordan Hall.}
\pRow{This agrees with what you said.}
\pRow{}
\pRow{(4) You said: Social work is a profession that is based in the philosophical tradition of humanism. It is an intellectual discipline that has its roots in the 1800s.}
\pRow{I checked: When did social work have its roots?}
\pRow{I found this article: The Emergence and Growth of the Social work Profession<br><br> Social work's roots were planted in the 1880s, when charity organization societies (COS) were created to organize municipal voluntary relief associations and settlement houses were established.}
\pRow{Social work has its roots in the 1880s, not 1800s.}
\pRow{This disagrees with what you said.}
\pRow{}
\pRow{(5) You said: The Havel-Hakimi algorithm is an algorithm for converting the adjacency matrix of a graph into its adjacency list. It is named after Vaclav Havel and Samih Hakimi.}
\pRow{I checked: What is the Havel-Hakimi algorithm?}
\pRow{I found this article: The Havel-Hakimi algorithm constructs a special solution if a simple graph for the given degree sequence exists, or proves that one cannot find a positive answer. This construction is based on a recursive algorithm. The algorithm was published by Havel (1955), and later by Hakimi (1962).}
\pRow{Havel-Hakimi algorithm is for constructing a special solution if a simple graph for the given degree sequence exists, or proving that one cannot find a positive answer, not converting the adjacency matrix of a graph into its adjacency list.}
\pRow{This disagrees with what you said.}
\pRow{}
\pRow{(6) You said: "Time of My Life" is a song by American singer-songwriter Bill Medley from the soundtrack of the 1987 film Dirty Dancing. The song was produced by Michael Lloyd.}
\pRow{I checked: Who was the producer of "(I've Had) The Time of My Life"?}
\pRow{I found this article: On September 8, 2010, the original demo of this song, along with a remix by producer Michael Lloyd, was released as digital files in an effort to raise money for the Patrick Swayze Pancreas Cancer Resarch Foundation at Stanford University.}
\pRow{"Time of My Life" was produced by Michael Lloyd.}
\pRow{This agrees with what you said.}
\pRow{}
\pRow{(7) You said: Kelvin Hopins was suspended from the Labor Party because he had allegedly sexually harassed and behaved inappropriately towards a Labour Party activist, Ava Etemadzadeh.}
\pRow{I checked: Why was Kelvin Hopins suspeneded from the Labor Party? }
\pRow{I found this article: A former Labour MP has left the party before an inquiry into sexual harassment allegations against him was able to be concluded, the party has confirmed. Kelvin Hopkins was accused in 2017 of inappropriate physical contact and was suspended by the Labour party pending an investigation.This agrees with what you said.}
\pRow{Kelvin Hopins was suspended because he had allegedly sexually harassed and behaved inappropriately towards a Labour Party activist, Ava Etemadzadeh.}
\pRow{This agrees with what you said.}
\pRow{}
\pRow{(8) You said: In the battles of Lexington and Concord, the British side was led by General Thomas Smith.}
\pRow{I checked: Who led the British side in the battle of Lexington and Concord?}
\pRow{I found this article: Interesting Facts about the Battles of Lexington and Concord. The British were led by Lieutenant Colonel Francis Smith. There were 700 British regulars.}
\pRow{The British side was led by Lieutenant Colonel Francis Smith, not General Thomas Hall.}
\pRow{This disagrees with what you said.}
\pRow{}
\pRow{(9) You said: \pVar{text}}
\pRow{I checked: \pVar{query}}
\pRow{I found this article: \pVar{evidence}}
\pRow{\pBlank}
&\\\hline\end{tabular}
\caption{Few-shot prompt for the agreement model, which uses chain-of-thought prompting.}
\label{fig:prompt_agreement}
\end{center}
\end{figure*}

\begin{figure*}[p]
\begin{center}\scriptsize
\begin{tabular}{|rp{\pWidth}|}\hline\pStart&\\
\pRow{[web] I will fix some things you said.}
\pRow{}
\pRow{(1) You said: Your nose switches back and forth between nostrils. When you sleep, you switch about every 45 minutes.  This is to prevent a buildup of mucus. It's called the nasal cycle.}
\pRow{I checked: How often do your nostrils switch?}
\pRow{I found this article: Although we don't usually notice it, during the nasal cycle one nostril becomes congested and thus contributes less to airflow, while the other becomes decongested. On average, the congestion pattern switches about every 2 hours, according to a small 2016 study published in the journal PLOS One.}
\pRow{This suggests 45 minutes switch time in your statement is wrong.}
\pRow{My fix: Your nose switches back and forth between nostrils. When you sleep, you switch about every 2 hours. This is to prevent a buildup of mucus. It's called the nasal cycle.}
\pRow{}
\pRow{(2) You said: In the battles of Lexington and Concord, the British side was led by General Thomas Hall.}
\pRow{I checked: who led the British side in the battle of Lexington and Concord?}
\pRow{I found this article: Interesting Facts about the Battles of Lexington and Concord. The British were led by Lieutenant Colonel Francis Smith. There were 700 British regulars.}
\pRow{This suggests General Thomas Hall in your statement is wrong.}
\pRow{My fix: In the battles of Lexington and Concord, the British side was led by Lieutenant Colonel Francis Smith.}
\pRow{}
\pRow{(3) You said: The Stanford Prison Experiment was conducted in the basement of Encina Hall, Stanford's psychology building.}
\pRow{I checked: where was Stanford Prison Experiment conducted.}
\pRow{I found this article: Carried out August 15-21, 1971 in the basement of Jordan Hall, the Stanford Prison Experiment set out to examine the psychological effects of authority and powerlessness in a prison environment.}
\pRow{This suggests Encina Hall in your statement is wrong.}
\pRow{My fix: The Stanford Prison Experiment was conducted in the basement of Jordan Hall, Stanford's psychology building.}
\pRow{}
\pRow{(4) You said: Phoenix Mills Ltd., a diversified business conglomerate, was established in 1854. It has a history of over 160 years.}
\pRow{I checked: When was Phoenix Mills Ltd. founded?}
\pRow{I found this article: Phoenix Mills Ltd was incorporated in the year 1905. The company began their operations as a textile manufacturing company on 17.3 acres of land at Lower Parel in Mumbai. In the year 1959 the company was listed in the Bombay Stock Exchange.}
\pRow{This suggests the year of establishment 1854 in your statement is wrong.}
\pRow{My fix: Phoenix Mills Ltd., a diversified business conglomerate, was established in 1905. It has a history of over 160 years.}
\pRow{}
\pRow{(5) You said: The Havel-Hakimi algorithm is an algorithm for converting the adjacency matrix of a graph into its adjacency list. It is named after Vaclav Havel and Samih Hakimi.}
\pRow{I checked: What is the Havel-Hakimi algorithm?}
\pRow{I found this article: The Havel-Hakimi algorithm constructs a special solution if a simple graph for the given degree sequence exists, or proves that one cannot find a positive answer. This construction is based on a recursive algorithm. The algorithm was published by Havel (1955), and later by Hakimi (1962).}
\pRow{This suggests the Havel-Hakimi algorithm's functionality in your statement is wrong.}
\pRow{My fix: The Havel-Hakimi algorithm constructs a special solution if a simple graph for the given degree sequence exists, or proves that one cannot find a positive answer. It is named after Vaclav Havel and Samih Hakimi}
\pRow{}
\pRow{(6) You said: "Time of My Life" is a song by American singer-songwriter Bill Medley from the soundtrack of the 1987 film Dirty Dancing. The song was produced by Phil Ramone.}
\pRow{I checked: Who was the producer of "(I've Had) The Time of My Life"?}
\pRow{I found this article: On September 8, 2010, the original demo of this song, along with a remix by producer Michael Lloyd, was released as digital files in an effort to raise money for the Patrick Swayze Pancreas Cancer Resarch Foundation at Stanford University.}
\pRow{This suggests "Time of My Life" producer name in your statement is wrong.}
\pRow{My fix: "Time of My Life" is a song by American singer-songwriter Bill Medley from the soundtrack of the 1987 film Dirty Dancing. The song was produced by Michael Lloyd.}
\pRow{}
\pRow{(7) You said: Phoenix Market City Pune is located on 21 acres of prime property in Pune. It is spread across four levels with approximately 1.4 million square feet of built-up space. The mall is owned and operated by Phoenix Mills Limited.}
\pRow{I checked: What is the area of Phoenix Market City in Pune?}
\pRow{I found this article: Phoenix Market City was opened in January 2013 and has the distinction of being the largest mall in the city of Pune, with the area of 3.4 million square feet. It is located in the Viman Nagar area of Pune.}
\pRow{This suggests the 1.4 million square feet of built-up space in your statment is wrong.}
\pRow{My fix: Phoenix Market City Pune is located on 21 acres of prime property in Pune. It is spread across four levels with approximately 3.4 million square feet of built-up space. The mall is owned and operated by Phoenix Mills Limited.}
\pRow{}
\pRow{(8) You said: \pVar{text}}
\pRow{I checked: \pVar{query}}
\pRow{I found this article: \pVar{evidence}}
\pRow{This suggests \pBlank}
&\\\hline\end{tabular}
\caption{Few-shot prompt for the revision model, which uses chain-of-thought prompting.}
\label{fig:prompt_revision}
\end{center}
\end{figure*}

\begin{figure*}[p]
\begin{center}\scriptsize
\begin{tabular}{|rp{\pWidth}|}\hline\pStart&\\
\pRow{[web] I will read the context and check only the last thing you said by asking questions.}
\pRow{}
\pRow{(1) Context: Your nose switches back and forth between nostrils. When you sleep, you switch about every 45 minutes.}
\pRow{You said: This is to prevent a buildup of mucus. It's called the nasal cycle.}
\pRow{To verify what you just said,}
\pRow{a) I googled: Why does your nostril switch during sleep?}
\pRow{b) I googled: What is nasal cycle?}
\pRow{c) I googled: What is the nostril switching during sleep called?}
\pRow{}
\pRow{(2) Context: The Stanford Prison Experiment was conducted in the basement of Encina Hall, Stanford's psychology building.}
\pRow{You said: It is a psychological study to observe the behaviors of conflict and violence that happen between inmates and prisoners in real prisons.}
\pRow{To verify what you just said,}
\pRow{a) I googled: What type of experiment was the Stanford Prison Experiment?}
\pRow{b) I googled: What was the objective of the Stanford Prison Experiment?}
\pRow{}
\pRow{(3) Context: The Havel-Hakimi algorithm is an algorithm for converting the adjacency matrix of a graph into its adjacency list.}
\pRow{You said: It is named after Vaclav Havel and Samih Hakimi.}
\pRow{To verify what you just said,}
\pRow{a) I googled: Who are Havel-Hakimi algorithm named after?}
\pRow{}
\pRow{(4) Context: "Time of My Life" is a song by American singer-songwriter Bill Medley from the soundtrack of the 1987 film Dirty Dancing.}
\pRow{You said: The song was produced by Michael Lloyd in the same year.}
\pRow{To verify what you just said,}
\pRow{a) I googled: Who produced the song "Time of My Life"?}
\pRow{b) I googled: When was the song "Time of My Life" by Bill Medley produced?}
\pRow{}
\pRow{(5) Context: The Late Show with Stephen Colbert is an American late-night talk show hosted by Stephen Colbert, which premiered on September 8, 2015.}
\pRow{You said: Produced by Spartina Productions and CBS Television Studios, it is the second iteration of CBS' Late Show franchise.}
\pRow{To verify what you just said,}
\pRow{a) I googled: Who produces "The Late Show with Stephen Colbert"?}
\pRow{b) I googled: What are the iterations of CBS' Late Show franchise?}
\pRow{}
\pRow{(6) Context: Super Mario Sunshine was released on GameCube in 2002. In the game, Mario uses a tool strapped to his back called FLUDD, which stands for The Flash Liquidizer Ultra Dousing Device.}
\pRow{You said: It can be used to spray water at objects or enemies. This allows Mario to change his movements, kill enemies, or clean up hazards on the floor.}
\pRow{To verify what you just said,}
\pRow{a) I googled: What is the main function of FLUDD in Super Mario Sunshine?}
\pRow{b) I googled: What can FLUDD in Super Mario Sunshine be used on?}
\pRow{c) I googled: In Super Mario Sunshine, can Mario change movement with FLUDD?}
\pRow{d) I googled: In Super Mario Sunshine, can Mario kill enemies with FLUDD?}
\pRow{e) I googled: In Super Mario Sunshine, can Mario clean up hazards on the floor with FLUDD?}
\pRow{}
\pRow{(7) Context: \pVar{context}}
\pRow{You said: \pVar{text}}
\pRow{To verify what you just said,}
\pRow{\pBlank}
&\\\hline\end{tabular}
\caption{Contextual version of the query generation prompt. The prompt works well for dialog contexts from QReCC even though the few-shot examples are not formatted as such.}
\label{fig:prompt_cqgen_contextual}
\end{center}
\end{figure*}

\begin{figure*}[p]
\begin{center}\scriptsize
\begin{tabular}{|rp{\pWidth}|}\hline\pStart&\\
\pRow{[web] I will check some things you said.}
\pRow{}
\pRow{(1) Context: Your nose switches back and forth between nostrils. It's called the nasal cycle. This is to prevent a buildup of mucus.}
\pRow{You said: When you sleep, you switch about every 45 minutes.}
\pRow{I checked: How often do your nostrils switch?}
\pRow{I found this article: Although we don't usually notice it, during the nasal cycle one nostril becomes congested and thus contributes less to airflow, while the other becomes decongested. On average, the congestion pattern switches about every 2 hours, according to a small 2016 study published in the journal PLOS One.}
\pRow{Your nose's switching time is about every 2 hours, not 45 minutes.}
\pRow{This disagrees with what you said.}
\pRow{}
\pRow{(2) Context: The Little House books is a series of American children's novels.}
\pRow{You said: The books were published by HarperCollins.}
\pRow{I checked: Who published the Little House books?}
\pRow{I found this article: These are the books that started it all -- the stories that captured the hearts and imaginations of children and young adults worldwide. Written by Laura Ingalls Wilder and published by HarperCollins, these beloved books remain a favorite to this day.}
\pRow{The Little House books were published by HarperCollins.}
\pRow{This agrees with what you said.}
\pRow{}
\pRow{(3) Context: The Stanford Prison Experiment is a psychological study to observe the behaviors of conflict and violence that happen between inmates and prisoners in real prisons.}
\pRow{You said: It was conducted in the basement of Jordan Hall, Stanford's psychology building.}
\pRow{I checked: Where was Stanford Prison Experiment conducted?}
\pRow{I found this article: Carried out August 15-21, 1971 in the basement of Jordan Hall, the Stanford Prison Experiment set out to examine the psychological effects of authority and powerlessness in a prison environment.}
\pRow{The Stanford Prison Experiment was conducted in Jordan Hall.}
\pRow{This agrees with what you said.}
\pRow{}
\pRow{(4) Context: Social work is a profession that is based in the philosophical tradition of humanism.}
\pRow{You said: It is an intellectual discipline that has its roots in the 1800s.}
\pRow{I checked: When did social work have its roots?}
\pRow{I found this article: The Emergence and Growth of the Social work Profession<br><br> Social work's roots were planted in the 1880s, when charity organization societies (COS) were created to organize municipal voluntary relief associations and settlement houses were established.}
\pRow{Social work has its roots in the 1880s, not 1800s.}
\pRow{This disagrees with what you said.}
\pRow{}
\pRow{(5) Context: The Havel-Hakimi algorithm is named after Vaclav Havel and Samih Hakimi.}
\pRow{You said: It is an algorithm for converting the adjacency matrix of a graph into its adjacency list.}
\pRow{I checked: What is the Havel-Hakimi algorithm?}
\pRow{I found this article: The Havel-Hakimi algorithm constructs a special solution if a simple graph for the given degree sequence exists, or proves that one cannot find a positive answer. This construction is based on a recursive algorithm. The algorithm was published by Havel (1955), and later by Hakimi (1962).}
\pRow{Havel-Hakimi algorithm is for constructing a special solution if a simple graph for the given degree sequence exists, or proving that one cannot find a positive answer, not converting the adjacency matrix of a graph into its adjacency list.}
\pRow{This disagrees with what you said.}
\pRow{}
\pRow{(6) Context: "Time of My Life" is a song by American singer-songwriter Bill Medley from the soundtrack of the 1987 film Dirty Dancing. }
\pRow{You said: The song was produced by Michael Lloyd in the same year.}
\pRow{I checked: Who was the producer of "(I've Had) The Time of My Life"?}
\pRow{I found this article: On September 8, 2010, the original demo of this song, along with a remix by producer Michael Lloyd, was released as digital files in an effort to raise money for the Patrick Swayze Pancreas Cancer Resarch Foundation at Stanford University.}
\pRow{The song "Time of My Life" was produced by Michael Lloyd.}
\pRow{This agrees with what you said.}
\pRow{}
\pRow{(7) Context: Super Mario Sunshine was released on GameCube in 2002. In the game, Mario uses a tool strapped to his back called FLUDD.}
\pRow{You said: FLUDD stands for Functional Language in a Unified Design Discipline. It can be used to spray water at objects or enemies. This allows Mario to change his movements, kill enemies, or clean up hazards on the floor.}
\pRow{I checked: What does FLUDD stands for in Super Mario Sunshine?}
\pRow{I found this article: The Flash Liquidizer Ultra Dousing Device, abbreviated and better known as FLUDD or F.L.U.D.D., is a multipurpose water pack from Super Mario Sunshine invented by Professor Elvin Gadd, indicated by the Gadd Science, Incorporated logo at the base of its nozzle exclusively during the cutscene at Pinna Park.}
\pRow{In Super Mario Sunshine, FLUDD stands for the Flash Liquidizer Ultra Dousing Device, not Functional Language in a Unified Design Discipline.}
\pRow{This disagrees with what you said.}
\pRow{}
\pRow{(8) Context: \pVar{context}}
\pRow{You said: \pVar{text}}
\pRow{I checked: \pVar{query}}
\pRow{I found this article: \pVar{evidence}}
\pRow{\pBlank}
&\\\hline\end{tabular}
\caption{Contextual version of the agreement model prompt.}
\label{fig:prompt_agreement_contextual}
\end{center}
\end{figure*}

\begin{figure*}[p]
\begin{center}\scriptsize
\begin{tabular}{|rp{\pWidth}|}\hline\pStart&\\
\pRow{[web] I will fix some things you said.}
\pRow{}
\pRow{(1) Context: Your nose switches back and forth between nostrils. It's called the nasal cycle. This is to prevent a buildup of mucus.}
\pRow{You said: When you sleep, you switch about every 45 minutes.}
\pRow{I checked: How often do your nostrils switch?}
\pRow{I found this article: Although we don't usually notice it, during the nasal cycle one nostril becomes congested and thus contributes less to airflow, while the other becomes decongested. On average, the congestion pattern switches about every 2 hours, according to a small 2016 study published in the journal PLOS One.}
\pRow{This suggests 45 minutes switch time in your statement is wrong.}
\pRow{My fix: When you sleep, you switch about every 2 hours.}
\pRow{}
\pRow{(2) Context: The Little House books is a series of American children's novels.}
\pRow{You said: The books were published by Amberjack Publishing.}
\pRow{I checked: Who published the Little House books?}
\pRow{I found this article: These are the books that started it all -- the stories that captured the hearts and imaginations of children and young adults worldwide. Written by Laura Ingalls Wilder and published by HarperCollins, these beloved books remain a favorite to this day.}
\pRow{This suggests Amberjack Publishing in your statement is wrong.}
\pRow{My fix: The books were published by HarperCollins.}
\pRow{}
\pRow{(3) Context: The Stanford Prison Experiment is a psychological study to observe the behaviors of conflict and violence that happen between inmates and prisoners in real prisons.}
\pRow{You said: It was conducted in the basement of Encina Hall, Stanford's psychology building.}
\pRow{I checked: where was Stanford Prison Experiment conducted.}
\pRow{I found this article: Carried out August 15-21, 1971 in the basement of Jordan Hall, the Stanford Prison Experiment set out to examine the psychological effects of authority and powerlessness in a prison environment.}
\pRow{This suggests Encina Hall in your statement is wrong.}
\pRow{My fix: It was conducted in the basement of Jordan Hall, Stanford's psychology building.}
\pRow{}
\pRow{(4) Context: The Havel-Hakimi algorithm is named after Vaclav Havel and Samih Hakimi.}
\pRow{You said: It is an algorithm for converting the adjacency matrix of a graph into its adjacency list.}
\pRow{I checked: What is the Havel-Hakimi algorithm?}
\pRow{I found this article: The Havel-Hakimi algorithm constructs a special solution if a simple graph for the given degree sequence exists, or proves that one cannot find a positive answer. This construction is based on a recursive algorithm. The algorithm was published by Havel (1955), and later by Hakimi (1962).}
\pRow{This suggests the Havel-Hakimi algorithm's functionality in your statement is wrong.}
\pRow{My fix: It constructs a special solution if a simple graph for the given degree sequence exists, or proves that one cannot find a positive answer.}
\pRow{}
\pRow{(5) Context: "Time of My Life" is a song by American singer-songwriter Bill Medley from the soundtrack of the 1987 film Dirty Dancing.}
\pRow{You said: The song was produced by Phil Ramone in the same year.}
\pRow{I checked: Who was the producer of "(I've Had) The Time of My Life"?}
\pRow{I found this article: On September 8, 2010, the original demo of this song, along with a remix by producer Michael Lloyd, was released as digital files in an effort to raise money for the Patrick Swayze Pancreas Cancer Resarch Foundation at Stanford University.}
\pRow{This suggests "Time of My Life" producer name in your statement is wrong.}
\pRow{My fix: The song was produced by Michael Lloyd in the same year.}
\pRow{}
\pRow{(6) Context: Phoenix Market City Pune is located on 21 acres of prime property in Pune. The mall is owned and operated by Phoenix Mills Limited.}
\pRow{You said: It is spread across four levels with approximately 1.4 million square feet of built-up space.}
\pRow{I checked: What is the area of Phoenix Market City in Pune?}
\pRow{I found this article: Phoenix Market City was opened in January 2013 and has the distinction of being the largest mall in the city of Pune, with the area of 3.4 million square feet. It is located in the Viman Nagar area of Pune.}
\pRow{This suggests the 1.4 million square feet of built-up space in your statment is wrong.}
\pRow{My fix: It is spread across four levels with approximately 3.4 million square feet of built-up space.}
\pRow{}
\pRow{(7) Context: \pVar{context}}
\pRow{You said: \pVar{text}}
\pRow{I checked: \pVar{query}}
\pRow{I found this article: \pVar{evidence}}
\pRow{This suggests \pBlank}
&\\\hline\end{tabular}
\caption{Contextual version of the revision model prompt.}
\label{fig:prompt_revision_contextual}
\end{center}
\end{figure*}

\begin{figure*}[p]
\begin{center}\scriptsize
\begin{tabular}{|rp{\pWidth}|}\hline\pStart&\\
\pRow{[web] I will think step by step and answer your question.}
\pRow{}
\pRow{Question: is growing seedless cucumber good for a gardener with entomophobia}
\pRow{Explanation: Entomophobia is a fear of insects. Plants need insects to pollinate them. Seedless fruits such as seedless cucumbers do not require pollination, so seedless fruits do not require insects. This makes good for people with entomophobia. }
\pRow{Answer: Yes}
\pRow{}
\pRow{Question: Who was british pm and viceroy during quit india movement?}
\pRow{Explanation: The Quit India Movement was launched in 8th August 1942. The british PM at that time was Winston Churchill. The british viceroy during the movement is Victor Hope, usually referred to as Lord Linlithgow. }
\pRow{Answer: Winston Churchil and Lord Linlithgow}
\pRow{}
\pRow{Question: Which year does game over man come out on netflix?}
\pRow{Explanation: Game Over, Man! is an action-comedy movie released in March 23, 2018, on Netflix. It's director is Kyle Newacheck. The movie stars Anders Holm, Adam DeVine, and Blake Anderson. }
\pRow{Answer: March 23, 2018}
\pRow{}
\pRow{Question: would it be very difficult for Nuno Gomes to dive to the Red Sea's deepest point?}
\pRow{Explanation: Nuno Gomes' deepest dive in the Red Sea to date is 317 metres. The Red Sea has a maximum depth of over 3,000 metres. So it would be difficult for Nuno Gomes to the deepest point of the Red sea.}
\pRow{Answer: Yes}
\pRow{}
\pRow{Question: Are chinchillas cold-blooded?}
\pRow{Explanation: Chinchillas are rodents, which are mammals. All mammals are warm-blooded.}
\pRow{Answer: No}
\pRow{}
\pRow{Question: what are the chemo drugs for gastric cancer?}
\pRow{Explanation: Chemotherapy can be given for gastric cancer before or after surgery. Usually you have a combination of 2 or 3 drugs. The most common types are: epirubicin, cisplatin and capecitabine (ECX) ; epirubicin, cisplatin and fluorouracil (ECF); or fluorouracil, folinic acid, oxaliplatin and docetaxel (FLOT).}
\pRow{Answer: epirubicin, cisplatin and capecitabine (ECX) ; epirubicin, cisplatin and fluorouracil (ECF); or fluorouracil, folinic acid, oxaliplatin and docetaxel (FLOT)}
\pRow{}
\pRow{Question: How many sacks does Clay Matthews have in his career?}
\pRow{Explanation: Clay Matthews III had 383 solo tackles, 136 assists, and 91.5 sacks. His father, Clay Matthews Jr., also palyed in NFL and had 69.5 sacks in his career. }
\pRow{Answer: 91.5 (Clay Matthews III) and 69.5 (Clay Matthews Jr.)}
\pRow{}
\pRow{Question: In the U.S., the most deaths are caused by}
\pRow{Explanation: The leading cause of death in the U.S. is heart disease, causing nearly 700,000 deaths annually. The second most common cause of deaths is cancer.}
\pRow{Answer: heart disease}
\pRow{}
\pRow{Question: the rank of indian economy in terms of nominal gdp is}
\pRow{Explanation: India's GDP is economy is \$2.94 trillion. It is the fifth-largest in the world. The top GDP contries are United States, China, Japan, Germany and India. }
\pRow{Answer: 5}
\pRow{}
\pRow{Question: \pVar{question}}
\pRow{Explanation: \pBlank}
&\\\hline\end{tabular}
\caption{The PaLM prompt for generating long-form answers to questions from NQ and SQA.}
\label{fig:prompt_passage_palm}
\end{center}
\end{figure*}

\begin{figure*}[p]
\begin{center}\scriptsize
\begin{tabular}{|rp{\pWidth}|}\hline\pStart&\\
\pRow{I will think step by step and answer your question.}
\pRow{}
\pRow{1. Question: Is growing seedless cucumber good for a gardener with entomophobia?}
\pRow{2. Explanation: Entomophobia is a fear of insects. Plants need insects to pollinate them. Seedless fruits such as seedless cucumbers do not require pollination so seedless fruits do not require insects. This is good for people with entomophobia.}
\pRow{3. Answer: Yes.}
\pRow{}
\pRow{1. Question: Who was British PM and Biceroy during Quit India Movement?}
\pRow{2. Explanation: The Quit India Movement was launched in 8th August 1942. The British PM at that time was Winston Churchill. The British Biceroy during the movement was Victor Hope, usually referred to as Lord Linlithgow.}
\pRow{3. Answer: Winston Churchil and Lord Linlithgow.}
\pRow{}
\pRow{1. Question: Which year does Game Over Man come out on Netflix?}
\pRow{2. Explanation: Game Over, Man! is an action-comedy movie. Its director is Kyle Newacheck. The movie stars Anders Holm, Adam DeVine, and Blake Anderson. The movie was released March 23, 2018 on Netflix.}
\pRow{3. Answer: March 23, 2018.}
\pRow{}
\pRow{1. Question: Would it be very difficult for Nuno Gomes to dive to the Red Sea's deepest point?}
\pRow{2. Explanation: Nuno Gomes' deepest dive in the Red Sea to date is 317 meters. The Red Sea has a maximum depth of over 3,000 meters. So it would be difficult for Nuno Gomes to dive to the deepest point of the Red Sea.}
\pRow{3. Answer: Yes.}
\pRow{}
\pRow{1. Question: Are chinchillas cold-blooded?}
\pRow{2. Explanation: Chinchillas are rodents. Rodents are mammals. All mammals are warm-blooded.}
\pRow{3. Answer: No.}
\pRow{}
\pRow{1. Question: What are the chemo drugs for gastric cancer?}
\pRow{2. Explanation: Chemotherapy can be given for gastric cancer before or after surgery. Usually you have a combination of 2 or 3 drugs. The most common types are: epirubicin, cisplatin and capecitabine (ECX) ; epirubicin, cisplatin and fluorouracil (ECF); or fluorouracil, folinic acid, oxaliplatin and docetaxel (FLOT).}
\pRow{3. Answer: Epirubicin, cisplatin and capecitabine (ECX) ; epirubicin, cisplatin and fluorouracil (ECF); or fluorouracil, folinic acid, oxaliplatin and docetaxel (FLOT).}
\pRow{}
\pRow{1. Question: How many sacks does Clay Matthews have in his career?}
\pRow{2. Explanation: Clay Matthews has been pro linebacker for 11 seasons. Ten of the seasons he played for Green Bay Packers, while his last season was with Los Angeles Rams, dedicated to his hometown. Clay Matthews III had 383 solo tackles, 136 assists, and 91.5 sacks. His father, Clay Matthews Jr., also played in NFL and had 69.5 sacks in his career.}
\pRow{3. Answer: 91.5 (Clay Matthews III) and 69.5 (Clay Matthews Jr.).}
\pRow{}
\pRow{1. Question: In the U.S., the most deaths are caused by?}
\pRow{2. Explanation: According to the Centers for Disease Control and Prevention, there were 3 million deaths in 2018. The leading cause of death in the U.S. is heart disease, causing nearly 700,000 deaths annually. The second most common cause of deaths is cancer.}
\pRow{3. Answer: Heart disease.}
\pRow{}
\pRow{1. Question: The rank of Indian economy in terms of nominal GDP is?}
\pRow{2. Explanation: The top GDP contries are United States, China, Japan, Germany, and India. India's GDP is economy is \$2.94 trillion. It is the fifth-largest economy in the world.}
\pRow{3. Answer: 5.}
\pRow{}
\pRow{1. Question: Who won season 16 on Dancing with the Stars?}
\pRow{2. Explanation: Dancing with the Stars is an American dance competition television series that premiered on June 1, 2005, on ABC. The 16th season premiered on March 18, 2013. It was won by Kellie Pickler.}
\pRow{3. Answer: Kellie Pickler.}
\pRow{}
\pRow{1. Question: \pVar{question}}
\pRow{2. Explanation: \pBlank}
&\\\hline\end{tabular}
\caption{The GPT-3 prompt for generating long-form answers to questions from NQ and SQA.}
\label{fig:prompt_passage_gpt3}
\end{center}
\end{figure*}

\begin{figure*}[p]
\begin{center}\scriptsize
\begin{tabular}{|rp{\pWidth}|}\hline\pStart&\\
\pRow{Hi, I will think step by step and answer your question.}
\pRow{Is growing seedless cucumber good for a gardener with entomophobia?}
\pRow{Yes. Entomophobia is a fear of insects. Plants need insects to pollinate them. Seedless cucumber fruit does not require pollination. So seedless cucumber does not require insects.}
\pRow{Here is another question. Which year does game over man come out on netflix?}
\pRow{Game Over, Man! is an American action-comedy movie. It was released on March 23, 2018, on Netflix. It is directed by Kyle Newacheck, written by Anders Holm, and starring himself, Adam DeVine, and Blake Anderson.}
\pRow{Here is another question. What are the chemo drugs for gastric cancer?}
\pRow{Chemotherapy can be given for gastric cancer before or after surgery. Usually you have a combination of 2 or 3 drugs. The most common types are: epirubicin, cisplatin and capecitabine (ECX) ; epirubicin, cisplatin and fluorouracil (ECF); or fluorouracil, folinic acid, oxaliplatin and docetaxel (FLOT).}
\pRow{Here is another question. Was N.W.A a album?}
\pRow{N.W.A was an American hip hop group from Compton, California.}
\pRow{Were they successful?}
\pRow{They has sold over 10 million units in the United States alone.}
\pRow{How many singles did it have?}
\pRow{The group NWA released 8 singles.}
\pRow{Here is another question. \pVar{$Q_1$}}
\pRow{\pVar{$A_1$}}
\pRow{...}
\pRow{\pVar{$Q_k$}}
\pRow{\pBlank}
&\\\hline\end{tabular}
\caption{The LaMDA prompt for generating answers to questions from QReCC. Each line is a conversation turn. The dialog context from QReCC contains rounds of questions and answers ($Q_1, A_1, Q_2, A_2, \dots, Q_k$).}
\label{fig:prompt_passage_lamda}
\end{center}
\end{figure*}

\begin{figure*}[p]
\begin{center}\scriptsize
\begin{tabular}{|rp{\pWidth}|}\hline\pStart&\\
\pRow{I will think step by step and answer your question.}
\pRow{}
\pRow{Is growing seedless cucumber good for a gardener with entomophobia?}
\pRow{Yes. Entomophobia is a fear of insects. Plants need insects to pollinate them. Seedless cucumber fruit does not require pollination. So seedless cucumber does not require insects.}
\pRow{}
\pRow{Which year does game over man come out on netflix?}
\pRow{Game Over, Man! is an American action-comedy movie. It was released on March 23, 2018, on Netflix. It is directed by Kyle Newacheck, written by Anders Holm, and starring himself, Adam DeVine, and Blake Anderson.}
\pRow{}
\pRow{What are the chemo drugs for gastric cancer?}
\pRow{Chemotherapy can be given for gastric cancer before or after surgery. Usually you have a combination of 2 or 3 drugs. The most common types are: epirubicin, cisplatin and capecitabine (ECX) ; epirubicin, cisplatin and fluorouracil (ECF); or fluorouracil, folinic acid, oxaliplatin and docetaxel (FLOT).}
\pRow{}
\pRow{Was N.W.A an album?}
\pRow{N.W.A was an American hip hop group from Compton, California.}
\pRow{Were they successful?}
\pRow{They has sold over 10 million units in the United States alone.}
\pRow{How many singles did they have?}
\pRow{N.W.A had eight singles, including "Straight Outta Compton", "Express Yourself", "Gangsta Gangsta", "Dopeman" and "Alwayz Into Somethin'".}
\pRow{}
\pRow{\pVar{$Q_1$}}
\pRow{\pVar{$A_1$}}
\pRow{...}
\pRow{\pVar{$Q_k$}}
\pRow{\pBlank}
&\\\hline\end{tabular}
\caption{The GPT-3 prompt for generating answers to questions from QReCC.}
\label{fig:prompt_passage_gpt3_conv}
\end{center}
\end{figure*}

} %
